\chapter{Обзор методов рекомендательных систем}

\section{Последовательные рекомендации: от коллаборативной фильтрации к трансформерам}

\subsection{Эволюция подходов к моделированию пользовательских предпочтений}

Задача рекомендательных систем состоит в предсказании предпочтений пользователя на основе его предыдущих взаимодействий с каталогом элементов. Исторически первыми подходами к решению этой задачи стали методы коллаборативной фильтрации~\cite{koren2009matrix}, основанные на предположении о том, что пользователи со схожими паттернами потребления будут одинаково оценивать новые элементы. В рамках данного семейства методов выделяются два направления: memory-based подходы, вычисляющие сходство непосредственно по матрице взаимодействий, и model-based подходы, строящие латентные представления пользователей и элементов.

Методы матричной факторизации, такие как SVD, ALS и BPR~\cite{rendle2009bpr}, стали де-факто стандартом model-based коллаборативной фильтрации на протяжении 2010-х годов. Они представляют каждого пользователя и каждый элемент в виде плотного вектора фиксированной размерности $d$, а предсказание строится как скалярное произведение соответствующих векторов:
\begin{equation}
    \hat{r}_{ui} = \mathbf{p}_u^\top \mathbf{q}_i,
\end{equation}
где $\mathbf{p}_u \in \mathbb{R}^d$~--- латентное представление пользователя $u$, а $\mathbf{q}_i \in \mathbb{R}^d$~--- латентное представление элемента~$i$. Данный подход обладает вычислительной эффективностью и хорошей масштабируемостью, однако не учитывает порядок взаимодействий, что существенно ограничивает его применимость в доменах с выраженной временно\'{й} динамикой предпочтений~--- таких как музыкальное потребление.

Осознание важности последовательного характера пользовательского поведения привело к формированию отдельного направления~--- последовательных рекомендаций (sequential recommendation). В отличие от классической коллаборативной фильтрации, данные методы рассматривают историю взаимодействий пользователя как упорядоченную последовательность $\mathcal{S}_u = (s_1, s_2, \ldots, s_t)$ и стремятся предсказать следующий элемент $s_{t+1}$, опираясь на динамику предшествующих действий.

Первые подходы к последовательным рекомендациям основывались на марковских цепях, моделирующих переходы между состояниями фиксированного порядка. Затем рекуррентные нейронные сети (RNN), в частности архитектуры GRU4Rec~\cite{hidasi2016session}, продемонстрировали способность улавливать более сложные долгосрочные зависимости в пользовательских сессиях. Свёрточные подходы, такие как Caser~\cite{tang2018caser}, предложили альтернативный взгляд, рассматривая последовательность взаимодействий как <<изображение>> и применяя свёрточные фильтры для извлечения паттернов различных масштабов.

Однако принципиальный прорыв в качестве последовательных рекомендаций произошёл с появлением архитектур на основе механизма самовнимания (self-attention), заимствованного из области обработки естественного языка.

\subsection{SASRec: механизм самовнимания для последовательных рекомендаций}

Модель Self-Attentive Sequential Recommendation (SASRec), предложенная Kang и McAuley~\cite{sasrec2018}, стала первой работой, успешно применившей механизм самовнимания к задаче последовательных рекомендаций. Авторы отмечают, что марковские модели эффективно захватывают локальные паттерны на основе нескольких последних взаимодействий, тогда как рекуррентные сети способны учитывать долгосрочные зависимости, однако при этом страдают от проблемы затухания градиентов и последовательного характера обучения. SASRec объединяет достоинства обоих подходов: механизм самовнимания позволяет модели адаптивно выбирать наиболее релевантные элементы из всей истории взаимодействий.

Архитектура SASRec основана на стеке блоков самовнимания. Для входной последовательности элементов $\mathcal{S} = (s_1, \ldots, s_n)$ модель формирует матрицу эмбеддингов $\mathbf{E} \in \mathbb{R}^{n \times d}$, к которой добавляются позиционные эмбеддинги $\mathbf{P} \in \mathbb{R}^{n \times d}$, кодирующие порядок элементов в последовательности:
\begin{equation}
    \hat{\mathbf{E}} = \mathbf{E} + \mathbf{P}.
\end{equation}

Каждый блок самовнимания вычисляет взвешенное представление через scaled dot-product attention:
\begin{equation}
    \mathrm{Attention}(\mathbf{Q}, \mathbf{K}, \mathbf{V}) = \mathrm{softmax}\!\left(\frac{\mathbf{Q}\mathbf{K}^\top}{\sqrt{d}}\right)\mathbf{V},
\end{equation}
где $\mathbf{Q}$, $\mathbf{K}$, $\mathbf{V}$~--- проекции входного представления. Для обеспечения авторегрессионного характера предсказания применяется каузальная маска, запрещающая модели <<заглядывать в будущее>>: элемент на позиции $i$ может <<видеть>> только элементы на позициях $j \leq i$.

SASRec продемонстрировала существенное превосходство над рекуррентными (GRU4Rec) и свёрточными (Caser) бейзлайнами на нескольких бенчмарках, при этом требуя на порядок меньше времени для обучения благодаря возможности параллельных вычислений. Статья получила более 3100 цитирований и заложила фундамент для всей последующей линейки трансформерных моделей рекомендаций.

\subsection{BERT4Rec: двунаправленное внимание и задача маскирования}

Развивая идеи SASRec, Sun и соавторы~\cite{bert4rec2019} предложили модель BERT4Rec, адаптирующую архитектуру BERT к задаче последовательных рекомендаций. Ключевая идея работы состоит в замене однонаправленного (каузального) предсказания на двунаправленное: модель обучается восстанавливать случайно замаскированные элементы последовательности, используя контекст как слева, так и справа от пропущенной позиции.

Формально, для входной последовательности $\mathcal{S} = (s_1, \ldots, s_n)$ случайная доля элементов $\rho$ заменяется специальным токеном \texttt{[MASK]}. Задача обучения~--- предсказать оригинальные идентификаторы замаскированных элементов, что соответствует задаче Cloze из области обработки естественного языка:
\begin{equation}
    \mathcal{L} = -\sum_{s_m \in \mathcal{S}_{\text{mask}}} \log P(s_m \mid \tilde{\mathcal{S}}),
\end{equation}
где $\tilde{\mathcal{S}}$~--- последовательность с замаскированными элементами, а $\mathcal{S}_{\text{mask}}$~--- множество замаскированных позиций.

Данный подход обеспечивает более полное использование контекста по сравнению с каузальной моделью SASRec, что особенно важно при моделировании пользовательского поведения: предпочтения на данный момент могут зависеть не только от предыдущих, но и от последующих действий в рамках одной сессии. На момент публикации BERT4Rec продемонстрировала улучшение по метрикам NDCG@10 и Hit Rate@10 в диапазоне 7--24\% относительно SASRec на четырёх бенчмарках (Amazon Beauty, Steam, MovieLens-1M, MovieLens-20M). Статья набрала более 2700 цитирований, и модель стала одним из наиболее распространённых бейзлайнов в области последовательных рекомендаций.

Вместе с тем, последующие исследования~\cite{petrov2022bert4rec_replicability} выявили проблемы воспроизводимости результатов BERT4Rec: при корректной настройке гиперпараметров SASRec оказывается конкурентоспособной или даже превосходит BERT4Rec на ряде датасетов. Это указывает на то, что преимущества двунаправленного внимания не являются безусловными и зависят от характеристик данных и протокола оценки.

\subsection{gSASRec: обобщённая функция потерь и масштабирование негативов}

Проблемы воспроизводимости и чувствительности к выбору функции потерь привели к появлению работы Petrov и Macdonald~\cite{gsasrec2023}, получившей награду Best Paper на конференции RecSys 2023. Авторы показывают, что значительная часть различий в производительности между SASRec и BERT4Rec обусловлена не архитектурными отличиями, а выбором функции потерь и количеством негативных примеров.

В работе предлагается обобщённая бинарная кросс-энтропия (gBCE), параметризованная количеством негативных примеров $t$:
\begin{equation}
    \mathcal{L}_{\mathrm{gBCE}} = -\log \sigma(\hat{y}_{pos}) - t \cdot \log(1 - \sigma(\hat{y}_{neg})),
\end{equation}
где $\hat{y}_{pos}$ и $\hat{y}_{neg}$~--- предсказанные оценки для положительного и негативного примеров соответственно, а $\sigma$~--- сигмоидная функция. Увеличение числа негативных примеров при обучении SASRec с gBCE позволяет достичь 9.47\% улучшения NDCG на датасете MovieLens-1M по сравнению с наилучшей конфигурацией BERT4Rec, при этом затрачивая на 73\% меньше времени обучения.

Данная работа имеет особую значимость для настоящего исследования по нескольким причинам. Во-первых, эксперименты на датасетах MovieLens-1M, Steam и Gowalla (более 1 миллиона элементов) подтверждают масштабируемость подхода, что соответствует масштабу рассматриваемой задачи. Во-вторых, Abbattista и соавторы~\cite{abbattista2024sequential} впоследствии применили gSASRec к датасету Yandex Music Event 2019 и показали, что дополнение модели персонализированными оценками популярности (personalized popularity scores) позволяет получить улучшение на 25--70\% в задаче музыкальных рекомендаций. Следует отметить, что в настоящей работе используется датасет YAMBDA~\cite{yambda2025}~--- существенно более крупный набор данных Яндекс.Музыки, опубликованный в 2025 году и содержащий 4.79 миллиарда взаимодействий от 1 миллиона пользователей с 9.4 миллионами треков, что на два порядка превышает объём Yandex Music Event 2019.

\subsection{Роль повторного потребления в музыкальных рекомендациях}

Музыкальный домен обладает специфической особенностью, отличающей его от рекомендаций фильмов, товаров или новостей: пользователи регулярно переслушивают одни и те же треки. Abbattista и соавторы~\cite{abbattista2024sequential} детально исследуют этот феномен на данных Yandex Music Event 2019 (более ранний и существенно меньший датасет Яндекс.Музыки по сравнению с YAMBDA) и показывают, что простой бейзлайн Personalized Most Popular~--- рекомендация наиболее часто прослушиваемых данным пользователем треков~--- оказывается неожиданно сильным конкурентом для нейросетевых моделей.

Авторы предлагают механизм Personalized Popularity Awareness (PPA), инъецирующий персонализированные оценки популярности в трансформерные модели SASRec, BERT4Rec и gSASRec. Оценка популярности для пары пользователь-трек вычисляется как нормализованная частота прослушиваний:
\begin{equation}
    \mathrm{pop}(u, i) = \frac{c(u, i)}{\max_{j \in \mathcal{I}_u} c(u, j)},
\end{equation}
где $c(u, i)$~--- количество прослушиваний трека $i$ пользователем $u$, а $\mathcal{I}_u$~--- множество всех треков, прослушанных пользователем. Данная оценка используется как дополнительная фича при формировании эмбеддингов элементов.

Эксперименты показывают, что учёт персонализированной популярности улучшает метрики всех трёх трансформерных моделей на 25--70\%, при этом наибольший эффект наблюдается именно на датасете Yandex Music, что объясняется высокой долей повторных прослушиваний в музыкальном потреблении. Данный результат указывает на необходимость явного моделирования паттернов повторного потребления при построении музыкальных рекомендательных систем~--- в том числе в контексте групповых рекомендаций.

Greer и соавторы~\cite{greer2025beyond} на данных Amazon Music дополнительно исследуют роль архитектурных компонентов трансформера в музыкальных рекомендациях. Авторы обнаруживают контринтуитивный результат: удаление каузальной маски и позиционных эмбеддингов из gSASRec приводит к улучшению F1 на 28\%. Это свидетельствует о том, что в музыкальном домене строгий хронологический порядок взаимодействий может быть менее информативен, чем в других доменах, и модель выигрывает от возможности свободно агрегировать информацию из всей истории прослушиваний.

\section{Модели машинного обучения в промышленных рекомендательных воронках}

\subsection{Многоэтапная архитектура рекомендательных систем}

Промышленные рекомендательные системы, обслуживающие каталоги из миллионов элементов, как правило, организованы в виде многоэтапной воронки (multi-stage pipeline). Такая архитектура обусловлена невозможностью применить тяжёлую модель ко всему каталогу в режиме реального времени и предполагает последовательное сужение множества кандидатов через несколько этапов с возрастающей сложностью моделей~\cite{covington2016youtube}.

Типичная воронка состоит из следующих этапов. На первом этапе (candidate generation) из каталога объёмом порядка $10^6$ элементов отбираются $10^3$--$10^4$ кандидатов с помощью быстрых методов: коллаборативной фильтрации, приближённого поиска ближайших соседей в пространстве эмбеддингов (ANN-индексы, такие как FAISS~\cite{johnson2019faiss}, ScaNN или HNSW), а также двухбашенных моделей (DSSM~\cite{huang2013dssm}, YouTube DNN~\cite{covington2016youtube}). На втором этапе (pre-ranking) лёгкая модель, как правило основанная на градиентном бустинге и ограниченном наборе признаков, сужает множество до $10^2$--$10^3$ элементов. На третьем этапе (ranking) тяжёлая модель с сотнями признаков формирует финальный ранжированный список. Такая архитектура используется в Яндекс.Музыке~\cite{volokitin2023yandex}, YouTube Music~\cite{google2024transformers}, Spotify и других крупных сервисах.

В контексте настоящей работы многоэтапная воронка имеет особое значение: различные модели машинного обучения применяются на разных этапах, и для каждого этапа существуют свои требования к скорости инференса, качеству ранжирования и набору используемых признаков.

\subsection{Градиентный бустинг: CatBoost как базовый ранкер}

Градиентный бустинг на решающих деревьях (GBDT) остаётся одним из наиболее широко применяемых семейств алгоритмов на этапах pre-ranking и ranking промышленных рекомендательных систем. Среди реализаций GBDT особое место занимает CatBoost, предложенный Prokhorenkova и соавторами~\cite{prokhorenkova2018catboost}. Алгоритм вносит два ключевых вклада в методологию градиентного бустинга.

Первый вклад~--- ordered boosting, решающий проблему утечки целевой переменной (target leakage), которая возникает при вычислении градиентов на тех же данных, которые использовались для построения предыдущих деревьев. CatBoost использует случайные перестановки обучающей выборки и вычисляет градиенты для каждого объекта только на основе моделей, обученных на предшествующих объектах в данной перестановке:
\begin{equation}
    g_i = \frac{\partial \mathcal{L}(y_i, F_{\sigma(i)-1}(x_i))}{\partial F_{\sigma(i)-1}(x_i)},
\end{equation}
где $\sigma$~--- случайная перестановка, а $F_{\sigma(i)-1}$~--- модель, обученная на объектах, предшествующих $i$ в перестановке $\sigma$.

Второй вклад~--- эффективная обработка категориальных признаков через target statistics с учётом упорядоченности объектов, что устраняет необходимость ручного кодирования и позволяет работать с признаками высокой кардинальности~--- такими как идентификаторы треков или пользователей.

В задаче музыкальных рекомендаций CatBoost естественным образом встраивается в этап pre-ranking и ranking: модель принимает на вход признаки пары (пользователь, трек)~--- статистики прослушиваний, временные паттерны, played ratio, длительность трека~--- и предсказывает вероятность положительного взаимодействия (лайка). Механизм \texttt{auto\_class\_weights} позволяет корректировать дисбаланс между классами, что критично при работе с имплицитной обратной связью, где положительные взаимодействия составляют малую долю от общего числа событий.

\subsection{DeepFM: объединение факторизационных машин и глубоких сетей}

В то время как градиентный бустинг эффективно работает с табличными признаками, он не моделирует взаимодействия признаков (feature interactions) напрямую~--- каждое дерево делает splits по одному признаку. Guo и соавторы~\cite{guo2017deepfm} предложили архитектуру DeepFM, объединяющую факторизационные машины (FM) и глубокую нейронную сеть в единой модели для предсказания вероятности клика (CTR prediction).

Архитектура DeepFM состоит из двух параллельных компонент с общим входным слоем:
\begin{equation}
    \hat{y} = \sigma\!\left(y_{\mathrm{FM}} + y_{\mathrm{DNN}}\right),
\end{equation}
где $y_{\mathrm{FM}}$ отвечает за моделирование взаимодействий признаков низкого порядка через факторизационную машину:
\begin{equation}
    y_{\mathrm{FM}} = w_0 + \sum_{i=1}^{n} w_i x_i + \sum_{i=1}^{n}\sum_{j=i+1}^{n} \langle \mathbf{v}_i, \mathbf{v}_j \rangle x_i x_j,
\end{equation}
а $y_{\mathrm{DNN}}$ моделирует взаимодействия высокого порядка через многослойный перцептрон. Ключевое преимущество архитектуры~--- общий слой эмбеддингов (shared embedding layer), который позволяет FM- и DNN-компонентам совместно обучать представления признаков без необходимости ручного feature engineering.

DeepFM получила более 2400 цитирований и стала одной из базовых архитектур для CTR-предсказания. В контексте рекомендательной воронки DeepFM и её развития~--- DCN v2~\cite{wang2021dcnv2}, DIN~\cite{zhou2018din}, DLRM~\cite{naumov2019dlrm}~--- являются кандидатами для этапа тяжёлого ранжирования, где модель может использовать сотни признаков и их попарные взаимодействия.

Сравнительный анализ показывает, что CatBoost и DeepFM дополняют друг друга: CatBoost обеспечивает высокую производительность на табличных данных с минимальной настройкой и хорошую интерпретируемость через feature importance, тогда как DeepFM лучше масштабируется при наличии разреженных категориальных признаков высокой кардинальности и способна автоматически обнаруживать сложные нелинейные взаимодействия. В промышленной практике часто используется комбинация обоих подходов: CatBoost на этапе pre-ranking для быстрой фильтрации, DeepFM или её аналоги~--- на этапе финального ранжирования.

\subsection{Behavior Sequence Transformer: учёт последовательности действий в ранкере}

Классические ранкеры (CatBoost, DeepFM) обрабатывают каждую пару (пользователь, элемент) независимо, используя агрегированные статистики поведения в качестве признаков. Однако порядок действий пользователя содержит важную информацию о его текущих предпочтениях. Chen и соавторы~\cite{chen2019bst} предложили Behavior Sequence Transformer (BST), интегрирующий трансформерный блок в архитектуру ранкера для моделирования последовательности поведения.

BST добавляет к стандартному набору признаков ранкера (user features, item features, context features) закодированное трансформером представление последовательности последних $n$ взаимодействий пользователя. Механизм self-attention позволяет модели определить, какие из предыдущих взаимодействий наиболее релевантны для текущего элемента-кандидата. Выходное представление трансформера конкатенируется с остальными признаками и подаётся в многослойный перцептрон для финального предсказания.

Данная архитектура была внедрена в Alibaba для ранжирования рекомендаций на Taobao и продемонстрировала статистически значимое улучшение CTR в A/B тестах. В контексте настоящей работы BST представляет интерес как пример гибридного подхода, объединяющего преимущества трансформерного кодирования последовательностей с возможностями классического ранкера по использованию произвольных признаков~--- включая групповые агрегации и контекстные сигналы.

\section{Большие языковые модели в рекомендательных системах}

\subsection{Предпосылки и таксономия подходов}

Успехи больших языковых моделей (Large Language Models, LLM) в задачах обработки естественного языка~--- генерации текста, ответов на вопросы, рассуждений~--- закономерно привлекли внимание исследователей в области рекомендательных систем. Wu и соавторы~\cite{wu2023survey_llm4rec} в обзорной работе, получившей около 500 цитирований, выделяют две основные парадигмы применения LLM к рекомендациям: дискриминативную (Discriminative LLM for Recommendation, DLLM4Rec) и генеративную (Generative LLM for Recommendation, GLLM4Rec).

В рамках дискриминативной парадигмы LLM используются для обогащения представлений элементов и пользователей: предобученные текстовые эмбеддинги служат дополнительными признаками в классических рекомендательных моделях. В рамках генеративной парадигмы задача рекомендации переформулируется как задача генерации текста: модель на входе получает текстовое описание предпочтений пользователя, а на выходе порождает идентификаторы или названия рекомендуемых элементов.

Помимо непосредственного применения в качестве рекомендателя, LLM могут выполнять вспомогательные роли. К ним относятся обогащение признаков (feature enrichment), то есть генерация текстовых описаний элементов для последующего получения эмбеддингов; инициализация эмбеддингов, то есть использование LLM-представлений в качестве начальных весов для коллаборативных моделей; а также объяснение рекомендаций, то есть генерация понятных пользователю обоснований выбора. Каждое из этих направлений имеет свою специфику, ограничения и перспективы, которые рассматриваются ниже.

\subsection{P5: рекомендации как обработка языка}

Работа Geng и соавторов~\cite{geng2022p5}, получившая около 660 цитирований, стала одной из первых крупных попыток объединить множество задач рекомендательных систем в единой text-to-text парадигме. Авторы предлагают модель P5 (Pretrain, Personalized Prompt, and Predict Paradigm), построенную на архитектуре T5, которая обрабатывает пять семейств задач через единый интерфейс естественного языка: предсказание рейтинга, последовательные рекомендации, объяснение, прямые рекомендации и обзорное суммирование.

Каждая задача формализуется как текстовый промпт. Например, для последовательной рекомендации входной промпт может иметь вид: <<User\_523 has interacted with items Item\_17, Item\_204, Item\_89 in order. Predict the next item>>, а модель генерирует идентификатор следующего элемента. Для обучения используются персонализированные промпты (personalized prompts), в которые встраиваются идентификаторы пользователей и элементов как специальные токены.

Принципиальное значение P5 состоит в демонстрации того, что единая языковая модель способна решать разнородные задачи рекомендаций без специализированных архитектур для каждой из них. Вместе с тем, модель имеет существенное ограничение: она работает с идентификаторами элементов, закодированными в словарь на этапе обучения, и не может рекомендовать элементы, отсутствующие в обучающей выборке. Это ограничивает применимость подхода в сценариях с быстро обновляющимся каталогом~--- характерных для музыкального домена.

\subsection{TALLRec: эффективное дообучение LLM для рекомендаций}

Bao и соавторы~\cite{bao2023tallrec} в работе TALLRec (Tuning framework for Aligning LLMs with Recommendation) адресуют ключевую практическую проблему: разрыв между задачами, на которых обучались LLM (генерация текста), и задачами рекомендаций (ранжирование элементов по предпочтениям). Авторы показывают, что прямое применение LLM через промпты (in-context learning) даёт неудовлетворительные результаты из-за недостаточного количества рекомендательных данных в предобучающем корпусе.

Для преодоления этого разрыва предлагается двухэтапная процедура дообучения модели LLaMA-7B. На первом этапе выполняется instruction tuning на небольшом объёме рекомендательных данных, сформулированных как инструкции: <<Given that the user has watched movies A, B, C and liked them, will the user like movie D? Answer: Yes/No>>. На втором этапе проводится дообучение на данных конкретного домена. Ключевой результат: модель демонстрирует значимое улучшение качества рекомендаций даже при обучении менее чем на 100 примерах, и при этом обладает кросс-доменной обобщающей способностью~--- модель, дообученная на рекомендациях фильмов, показывает улучшение и в домене книг.

Практическая ценность TALLRec для настоящей работы определяется двумя факторами. Во-первых, подход работает на одной GPU (RTX 3090), что совместимо с ограничениями Google Colab. Во-вторых, кросс-доменная обобщающая способность открывает возможность применения LLM, дообученного на текстовых данных, к музыкальному домену~--- даже при отсутствии текстовых метаданных о треках в исходном датасете.

\subsection{Оценка способности LLM понимать пользовательские предпочтения}

Критический вопрос при использовании LLM в рекомендациях~--- насколько хорошо языковые модели действительно понимают пользовательские предпочтения? Kang и соавторы~\cite{kang2024llm_preferences} проводят систематическую оценку LLM различных масштабов (от 250M до 540B параметров) в задаче предсказания пользовательских рейтингов.

Авторы формулируют задачу как текстовый промпт, содержащий историю оценок пользователя, и просят модель предсказать рейтинг для нового элемента. Эксперименты на данных MovieLens и Amazon Reviews приводят к нескольким важным выводам. Во-первых, zero-shot LLM систематически уступают даже простым коллаборативным бейзлайнам (matrix factorization), имеющим доступ к матрице взаимодействий. Во-вторых, производительность LLM растёт с увеличением масштаба модели, но даже крупнейшие модели (PaLM 540B) не достигают уровня специализированных рекомендательных систем. В-третьих, LLM демонстрируют значительный bias в сторону популярных элементов и испытывают затруднения при работе с нишевыми предпочтениями.

Данные результаты имеют принципиальное значение для выбора архитектуры настоящей работы: они аргументируют в пользу гибридного подхода, в котором LLM не заменяет, а дополняет классические рекомендательные модели. В контексте групповых музыкальных рекомендаций LLM может использоваться для <<понимания контекста>> группы и обогащения признаков, тогда как финальное ранжирование выполняется специализированной моделью (CatBoost, трансформер), обученной на данных взаимодействий.

\subsection{HLLM: иерархические языковые модели для рекомендаций}

Chen и соавторы~\cite{chen2024hllm} из ByteDance предлагают архитектуру Hierarchical Large Language Model (HLLM), адресующую вопрос о реальной ценности предобученных весов LLM для рекомендаций. Архитектура состоит из двух уровней: Item LLM, извлекающая богатые контентные представления из текстовых описаний элементов, и User LLM, моделирующая интересы пользователя на основе последовательности таких представлений.

Формально, для элемента $i$ с текстовым описанием $\mathbf{t}_i$ Item LLM формирует контентный эмбеддинг:
\begin{equation}
    \mathbf{h}_i = \mathrm{ItemLLM}(\mathbf{t}_i) \in \mathbb{R}^d,
\end{equation}
а User LLM обрабатывает последовательность контентных эмбеддингов взаимодействий пользователя для предсказания следующего элемента:
\begin{equation}
    \hat{\mathbf{h}}_{t+1} = \mathrm{UserLLM}(\mathbf{h}_{s_1}, \mathbf{h}_{s_2}, \ldots, \mathbf{h}_{s_t}).
\end{equation}

Ключевые экспериментальные результаты: HLLM превосходит традиционные ID-based модели с большим отрывом на датасетах PixelRec и Amazon Reviews, при этом демонстрируя хорошую масштабируемость~--- наибольшая конфигурация использует 7B параметров как для Item LLM, так и для User LLM. Модель была проверена в production A/B тестах на сервисах ByteDance, что подтверждает её практическую применимость.

Для настоящей работы HLLM представляет особый интерес как пример архитектуры, в которой LLM-компонента может компенсировать отсутствие метаданных о треках. В датасете YAMBDA~\cite{yambda2025} отсутствуют жанры и имена артистов в текстовом виде, однако доступны анонимизированные аудиоэмбеддинги треков; Item LLM потенциально может быть использована для обогащения этих представлений, а полученные эмбеддинги~--- в качестве признаков для этапа ранжирования.

\section{Групповые рекомендации}

\subsection{Постановка задачи и классические стратегии агрегации}

Задача групповых рекомендаций (group recommendation) состоит в формировании списка элементов, удовлетворяющего предпочтениям не одного пользователя, а группы из $m$ участников $\mathcal{G} = \{u_1, u_2, \ldots, u_m\}$. Данная задача возникает в практических сценариях, когда группа людей совместно потребляет контент: выбор фильма для семейного просмотра, составление плейлиста для вечеринки, подбор ресторана для компании. В контексте настоящей работы рассматривается сценарий, в котором группа из нескольких человек собирается вместе и нуждается в музыкальных рекомендациях, учитывающих вкусы всех участников.

Классические подходы к групповым рекомендациям основаны на стратегиях агрегации индивидуальных предпочтений. Пусть $r(u, i)$~--- предсказанная оценка элемента $i$ для пользователя $u$. Тогда групповая оценка $r(\mathcal{G}, i)$ вычисляется одной из следующих стратегий:

\begin{itemize}
    \item \textbf{Average (AVG)}~--- среднее предсказанных оценок:
    \[
        r(\mathcal{G}, i) = \frac{1}{m}\sum_{u \in \mathcal{G}} r(u, i).
    \]
    \item \textbf{Least Misery (LM)}~--- минимальная оценка среди участников:
    \[
        r(\mathcal{G}, i) = \min_{u \in \mathcal{G}} r(u, i).
    \]
    Эта стратегия гарантирует отсутствие элементов, крайне неприятных хотя бы одному участнику;
    \item \textbf{Maximum Pleasure (MP)}~--- максимальная оценка:
    \[
        r(\mathcal{G}, i) = \max_{u \in \mathcal{G}} r(u, i).
    \]
    \item \textbf{Average without Misery (AwM)}~--- среднее с исключением элементов, получивших оценку ниже порога хотя бы от одного участника;
    \item \textbf{Multiplicative}~--- произведение нормализованных оценок:
    \[
        r(\mathcal{G}, i) = \prod_{u \in \mathcal{G}} r(u, i).
    \]
\end{itemize}

Данные стратегии обладают простотой реализации и интерпретируемостью, однако имеют существенное ограничение: они не учитывают взаимодействия между участниками группы и рассматривают каждого пользователя изолированно. В реальности предпочтения людей в группе могут зависеть от социальной динамики~--- например, участник может быть готов к компромиссу ради комфорта остальных, или влияние отдельных участников может быть неравномерным.

\subsection{AGREE: attention-механизм для групповых рекомендаций}

Cao, He и соавторы~\cite{cao2018agree} предложили модель AGREE (Attentive Group Recommendation)~--- первую работу, применившую attention-механизм к задаче групповых рекомендаций. Модель получила более 300 цитирований и заложила основу для нейросетевого подхода к агрегации групповых предпочтений.

AGREE построена на фреймворке Neural Collaborative Filtering (NCF) и состоит из двух компонент. Первая~--- attention-модуль, обучающийся динамически взвешивать вклад каждого участника группы:
\begin{equation}
    \alpha_u = \frac{\exp(\mathbf{w}^\top \mathbf{e}_u)}{\sum_{u' \in \mathcal{G}} \exp(\mathbf{w}^\top \mathbf{e}_{u'})},
\end{equation}
где $\mathbf{e}_u$~--- латентное представление пользователя $u$, а $\mathbf{w}$~--- обучаемый вектор весов attention. Групповое представление формируется как взвешенная сумма представлений участников:
\begin{equation}
    \mathbf{g} = \sum_{u \in \mathcal{G}} \alpha_u \cdot \mathbf{e}_u.
\end{equation}

Вторая компонента~--- нейросетевой предиктор, вычисляющий оценку совместимости группового представления $\mathbf{g}$ с представлением элемента $\mathbf{e}_i$ через многослойный перцептрон.

Ключевое преимущество AGREE перед классическими стратегиями агрегации состоит в способности модели автоматически определять <<влиятельность>> каждого участника для конкретной группы. Кроме того, модель способна обрабатывать cold-start сценарии: для пользователей без индивидуальных взаимодействий attention-веса могут быть вычислены на основе их принадлежности к группам, в которых они участвовали ранее.

\subsection{GroupIM: максимизация взаимной информации для эфемерных групп}

Sankar и соавторы~\cite{sankar2020groupim} в работе GroupIM адресуют проблему, особенно актуальную для настоящего исследования: рекомендации для эфемерных (ephemeral) групп~--- групп, формирующихся спонтанно и не имеющих длительной истории совместных взаимодействий. Именно такой сценарий характерен для ситуации, когда несколько человек собираются вместе и хотят послушать музыку.

В отличие от AGREE, GroupIM регуляризует латентные представления пользователей и групп через принцип максимизации взаимной информации (mutual information maximization). Модель решает три задачи одновременно: предсказание групповых предпочтений, согласованность представлений пользователь-группа (user-group mutual information), и информативный отбор наиболее значимых участников для предсказания.

Формально, помимо стандартной функции потерь рекомендаций $\mathcal{L}_{\mathrm{rec}}$, модель оптимизирует дополнительный регуляризатор:
\begin{equation}
    \mathcal{L} = \mathcal{L}_{\mathrm{rec}} + \lambda \cdot \mathcal{L}_{\mathrm{MI}},
\end{equation}
где $\mathcal{L}_{\mathrm{MI}}$ максимизирует взаимную информацию между представлениями отдельных пользователей и групп, в которые они входят. Это обеспечивает два эффекта: представления пользователей обогащаются информацией о групповом поведении, а групповые представления~--- информацией о наиболее информативных участниках.

Эксперименты показывают улучшение 31--62\% по метрике NDCG@20 по сравнению с предшествующими методами, включая AGREE. Особую ценность для настоящей работы представляет способность GroupIM динамически приоритизировать наиболее информативных членов группы~--- это может быть полезно в сценарии, когда часть участников имеет богатую историю прослушиваний, а часть является <<холодными>> пользователями.

\subsection{AlignGroup: согласование группового консенсуса с индивидуальными предпочтениями}

Xu и соавторы~\cite{xu2024aligngroup} в работе AlignGroup, представленной на CIKM 2024, предлагают наиболее актуальный подход к нейросетевым групповым рекомендациям. Авторы выделяют ключевую проблему предшествующих методов: модели, оптимизированные на групповых взаимодействиях, часто теряют точность в моделировании индивидуальных предпочтений, и наоборот.

AlignGroup решает эту проблему через двухуровневую архитектуру. На первом уровне гиперграфовая нейронная сеть (hypergraph neural network) моделирует внутригрупповые (intra-group) и межгрупповые (inter-group) связи. Гиперграф позволяет естественным образом представить группу как гиперребро, соединяющее нескольких пользователей, и агрегировать информацию как внутри групп, так и между пользователями, принадлежащими к разным группам с пересекающимся составом.

На втором уровне вводится задача self-supervised alignment, согласующая групповой консенсус с индивидуальными предпочтениями. Интуитивно, групповое представление должно отражать <<общий знаменатель>> вкусов участников, но не должно противоречить индивидуальным предпочтениям каждого из них. Alignment loss минимизирует расхождение между предсказаниями модели для группы и агрегированными предсказаниями для отдельных участников:
\begin{equation}
    \mathcal{L}_{\mathrm{align}} = \sum_{\mathcal{G}} \sum_{i \in \mathcal{I}_{\mathcal{G}}} \left\| f(\mathbf{g}, \mathbf{e}_i) - \mathrm{Agg}\left(\{f(\mathbf{e}_u, \mathbf{e}_i)\}_{u \in \mathcal{G}}\right) \right\|^2,
\end{equation}
где $f(\cdot, \cdot)$~--- функция предсказания, а $\mathrm{Agg}$~--- функция агрегации.

AlignGroup превосходит все предшествующие модели как в групповых, так и в индивидуальных рекомендациях, что делает данный подход особенно перспективным для настоящей работы, где система должна одновременно учитывать групповой контекст и индивидуальные предпочтения каждого участника.

\subsection{Справедливость в групповых рекомендациях}

Важным аспектом групповых рекомендаций является справедливость (fairness)~--- обеспечение того, чтобы рекомендации не систематически игнорировали предпочтения отдельных участников группы. Jin и соавторы~\cite{jin2023fairness_survey} в обзорной работе, получившей более 200 цитирований, систематизируют подходы к fairness в рекомендательных системах.

В контексте групповых рекомендаций выделяются несколько типов справедливости. Individual fairness требует, чтобы похожие пользователи получали похожие рекомендации. Group fairness требует, чтобы определённые подгруппы (например, по демографическим признакам) не были дискриминированы. Для задачи агрегации групповых предпочтений наиболее релевантна концепция proportional fairness~--- каждый участник группы должен получить <<справедливую долю>> удовлетворённости от итогового списка рекомендаций.

Формально, для оценки справедливости группового списка $\mathcal{L}$ используются метрики:
\begin{equation}
    \mathrm{Fairness}(\mathcal{G}, \mathcal{L}) = 1 - \frac{\mathrm{Var}\left(\{\mathrm{sat}(u, \mathcal{L})\}_{u \in \mathcal{G}}\right)}{\mathrm{MaxVar}},
\end{equation}
где $\mathrm{sat}(u, \mathcal{L})$~--- удовлетворённость пользователя $u$ списком $\mathcal{L}$, а дисперсия удовлетворённости измеряет неравномерность распределения <<выгоды>> между участниками. Стратегия Least Misery оптимизирует fairness, однако может жертвовать общим качеством рекомендаций; fairness-aware подходы стремятся найти оптимальный баланс между общим качеством и равномерностью удовлетворённости.

Для настоящей работы метрики справедливости представляют особую важность: при формировании группового плейлиста необходимо не только максимизировать среднюю удовлетворённость, но и гарантировать, что ни один участник не будет систематически <<обделён>> рекомендациями, не соответствующими его вкусам.

\section{Проблема холодного старта, контентные признаки и новые архитектуры}

\subsection{Проблема холодного старта в рекомендательных системах}

Проблема холодного старта (cold start) является одной из наиболее фундаментальных проблем рекомендательных систем и возникает в двух основных сценариях: появление нового пользователя (user cold start), для которого отсутствует история взаимодействий, и появление нового элемента (item cold start), для которого нет накопленных оценок или событий потребления. Bei и соавторы~\cite{bei2025coldstart_survey} в обзорной работе, охватывающей 220 публикаций до конца 2024 года, систематизируют подходы к холодному старту по четырём <<областям знания>> (knowledge scopes): контентные признаки (content features), графовые связи (graph relations), кросс-доменная информация (domain information) и знания больших языковых моделей (LLM world knowledge).

В контексте групповых музыкальных рекомендаций проблема холодного старта имеет особую остроту. Во-первых, при формировании группы некоторые участники могут не иметь истории прослушиваний в системе. Во-вторых, каталог музыкальных сервисов постоянно пополняется новыми треками, для которых отсутствуют данные о взаимодействиях. В датасете YAMBDA~\cite{yambda2025} данная проблема частично компенсируется наличием аудиоэмбеддингов для каждого трека, что позволяет использовать контентные признаки даже при отсутствии коллаборативных сигналов.

\subsection{DropoutNet: обучение через имитацию холодного старта}

Volkovs, Yu и Poutanen~\cite{volkovs2017dropoutnet} предложили элегантный подход к проблеме холодного старта, основанный на наблюдении, что сценарий холодного старта эквивалентен задаче обработки отсутствующих данных (missing data). Работа, представленная на NeurIPS 2017 и получившая более 400 цитирований, вводит архитектуру DropoutNet.

Ключевая идея состоит в следующем: стандартные модели коллаборативной фильтрации обучаются на латентных представлениях пользователей и элементов, полученных из матрицы взаимодействий. Для <<тёплых>> пользователей и элементов эти представления информативны, но для новых~--- отсутствуют. DropoutNet решает эту проблему, объединяя два источника информации в единой нейронной сети:
\begin{equation}
    \hat{r}_{ui} = f_\theta\!\left(\mathrm{NN}_u(\mathbf{p}_u, \mathbf{x}_u),\; \mathrm{NN}_i(\mathbf{q}_i, \mathbf{x}_i)\right),
\end{equation}
где $\mathbf{p}_u$, $\mathbf{q}_i$~--- латентные представления из коллаборативной модели, а $\mathbf{x}_u$, $\mathbf{x}_i$~--- контентные признаки (демография пользователя, аудиопризнаки трека). Во время обучения к латентным представлениям применяется dropout с высокой вероятностью, что вынуждает сеть опираться на контентные признаки~--- имитируя ситуацию холодного старта. В результате обученная модель естественным образом работает как с <<тёплыми>>, так и с <<холодными>> пользователями и элементами.

Данный подход особенно релевантен для настоящей работы: в датасете YAMBDA доступны аудиоэмбеддинги треков, которые могут служить контентными признаками $\mathbf{x}_i$ для новых треков, а DropoutNet может быть интегрирован в этап pre-ranking рекомендательной воронки.

\subsection{Мета-обучение для быстрой адаптации: MeLU}

Альтернативный подход к холодному старту предлагает парадигма мета-обучения (meta-learning), направленная на быструю адаптацию модели к новым пользователям по минимальному количеству взаимодействий. Lee и соавторы~\cite{lee2019melu} в работе MeLU (Meta-Learned User Preference Estimator), представленной на KDD 2019, адаптируют алгоритм MAML (Model-Agnostic Meta-Learning) к задаче рекомендаций.

MAML обучает начальные параметры модели $\theta_0$ таким образом, чтобы несколько шагов градиентного спуска на данных нового пользователя $u$ приводили к хорошему предсказанию его предпочтений:
\begin{equation}
    \theta_u = \theta_0 - \alpha \nabla_\theta \mathcal{L}_u^{\mathrm{support}}(\theta_0),
\end{equation}
где $\mathcal{L}_u^{\mathrm{support}}$~--- функция потерь на <<опорном>> (support) множестве из нескольких взаимодействий нового пользователя, а $\alpha$~--- скорость адаптации. Мета-обучение оптимизирует $\theta_0$ по множеству задач (пользователей), минимизируя потери на <<целевом>> (query) множестве после адаптации:
\begin{equation}
    \theta_0^* = \arg\min_{\theta_0} \sum_{u \in \mathcal{U}} \mathcal{L}_u^{\mathrm{query}}(\theta_u).
\end{equation}

MeLU дополнительно предлагает стратегию отбора <<свидетельских>> (evidence) элементов~--- треков, которые наиболее эффективно раскрывают предпочтения нового пользователя. Эксперименты показывают снижение MAE не менее чем на 5.92\% по сравнению с базовыми подходами.

В контексте групповых рекомендаций мета-обучение открывает возможность быстрой адаптации к <<холодным>> участникам группы: достаточно нескольких оценок от нового пользователя, чтобы модель сформировала приемлемое представление о его предпочтениях и включила их в групповую агрегацию.

\subsection{Контентные признаки в музыкальных рекомендациях}

Музыкальный домен обладает уникальным преимуществом перед другими доменами рекомендаций: аудиосигнал трека несёт в себе богатую информацию о его характеристиках~--- жанре, темпе, настроении, инструментовке. Это позволяет строить контентные представления треков даже при полном отсутствии коллаборативных данных.

Van den Oord и соавторы~\cite{vandeoord2013deep_music} в одной из первых работ на стыке глубокого обучения и музыкальных рекомендаций, получившей более 1600 цитирований, предложили использовать свёрточные нейронные сети (CNN) для извлечения латентных представлений из спектрограмм аудиозаписей. Авторы показывают, что обученные на аудио представления коррелируют с латентными факторами коллаборативной фильтрации, и что модель на основе аудио способна предсказывать предпочтения пользователей для треков, отсутствующих в матрице взаимодействий.

Данный результат получил дальнейшее развитие в работах по мультимодальным рекомендациям. В контексте настоящей работы наличие предвычисленных аудиоэмбеддингов в датасете YAMBDA~\cite{yambda2025} предоставляет непосредственную возможность построения контентных представлений треков без необходимости обучения собственных аудиомоделей. Эти эмбеддинги могут использоваться в нескольких ролях: как признаки для CatBoost на этапе ранжирования, как инициализация элементных эмбеддингов в трансформерных моделях (SASRec, gSASRec), и как основа для рекомендации <<холодных>> треков через DropoutNet-подобные архитектуры.

\subsection{Mamba и модели пространства состояний для последовательных рекомендаций}

Трансформерные архитектуры, ставшие основой современных последовательных рекомендательных систем (SASRec, BERT4Rec, gSASRec), обладают фундаментальным ограничением: квадратичная вычислительная сложность механизма self-attention $O(n^2 d)$ относительно длины последовательности $n$. Для пользователей с длинной историей прослушиваний (в YAMBDA средний пользователь имеет тысячи взаимодействий) это создаёт существенные вычислительные ограничения как при обучении, так и при инференсе.

Модели пространства состояний (State Space Models, SSM) предлагают альтернативу, обеспечивающую линейную сложность $O(n)$ при сохранении способности моделировать долгосрочные зависимости. Gu и Dao~\cite{gu2023mamba} в работе Mamba, получившей более 4000 цитирований, предложили архитектуру селективных моделей пространства состояний (selective SSM, или S6). Ключевая инновация Mamba~--- механизм селекции, делающий параметры SSM зависимыми от входного токена:
\begin{equation}
    \mathbf{h}_t = \bar{\mathbf{A}}_t \mathbf{h}_{t-1} + \bar{\mathbf{B}}_t \mathbf{x}_t, \qquad y_t = \mathbf{C}_t \mathbf{h}_t,
\end{equation}
где $\bar{\mathbf{A}}_t$, $\bar{\mathbf{B}}_t$, $\mathbf{C}_t$~--- дискретизированные матрицы, параметризованные текущим входом $\mathbf{x}_t$ через линейные проекции. В отличие от классических SSM (S4), где параметры фиксированы и не зависят от содержания последовательности, Mamba позволяет модели селективно <<запоминать>> или <<забывать>> информацию в зависимости от текущего контекста~--- аналогично механизму гейтов в LSTM, но с существенно более эффективной реализацией через параллельный scan-алгоритм.

Liu и соавторы~\cite{liu2024mamba4rec} в работе Mamba4Rec, получившей награду Best Paper на RelKD Workshop (KDD 2024), стали первыми, кто применил архитектуру Mamba к задаче последовательных рекомендаций. Авторы адаптируют блоки Mamba для обработки последовательностей взаимодействий пользователя, добавляя позиционные эмбеддинги и слой нормализации для стабилизации обучения.

Эксперименты на бенчмарках Amazon Beauty, Sports, Toys и MovieLens-1M показывают, что Mamba4Rec превосходит как RNN-бейзлайны (GRU4Rec), так и трансформерные модели (SASRec, BERT4Rec) по метрикам HR@$k$ и NDCG@$k$, при этом обеспечивая значительно более быстрый инференс на длинных последовательностях благодаря линейной сложности и отсутствию необходимости хранить матрицу attention.

Для настоящей работы Mamba4Rec представляет интерес по двум причинам. Во-первых, линейная сложность позволяет обрабатывать полные истории прослушиваний пользователей YAMBDA без необходимости усечения последовательностей~--- в отличие от трансформерных моделей, обычно ограниченных последними 50--200 взаимодействиями. Во-вторых, архитектура Mamba4Rec совместима с вычислительными ограничениями Google Colab: авторы отмечают, что модель обучается быстрее SASRec при сопоставимом или меньшем потреблении памяти.

\section{Выводы по главе}

Проведённый обзор литературы позволяет сформулировать ряд выводов, определяющих архитектурные и методологические решения настоящей работы.

Во-первых, область последовательных рекомендаций прошла путь от матричной факторизации и рекуррентных сетей к трансформерным моделям (SASRec, BERT4Rec), при этом ключевым фактором качества оказывается не столько архитектура, сколько выбор функции потерь и стратегия негативного сэмплирования~\cite{gsasrec2023}. Обобщённая бинарная кросс-энтропия (gBCE) обеспечивает до 9.47\% улучшения NDCG при значительно меньших вычислительных затратах, что делает gSASRec предпочтительным бейзлайном для этапа генерации кандидатов. Специфика музыкального домена~--- высокая доля повторного потребления и меньшая роль хронологического порядка~\cite{greer2025beyond}~--- требует адаптации стандартных трансформерных моделей, в частности включения механизмов персонализированной популярности~\cite{abbattista2024sequential}.

Во-вторых, промышленные рекомендательные системы организованы в виде многоэтапной воронки, где на каждом этапе используются модели с различным соотношением скорости и качества. CatBoost~\cite{prokhorenkova2018catboost} естественным образом подходит для этапов pre-ranking и ranking благодаря встроенной обработке категориальных признаков и механизму ordered boosting, устраняющему утечку целевой переменной. DeepFM~\cite{guo2017deepfm} и BST~\cite{chen2019bst} демонстрируют преимущества нейросетевых ранкеров при наличии разреженных признаков высокой кардинальности и необходимости учёта последовательности поведения.

В-третьих, большие языковые модели открывают новые возможности для рекомендательных систем, однако в текущем состоянии не способны заменить специализированные модели: zero-shot LLM систематически уступают коллаборативным методам~\cite{kang2024llm_preferences}. Наиболее перспективным представляется гибридный подход, в котором LLM используются для обогащения признаков (HLLM~\cite{chen2024hllm}) или быстрой адаптации к домену (TALLRec~\cite{bao2023tallrec}), а финальное ранжирование выполняется классическими моделями.

В-четвёртых, задача групповых рекомендаций требует перехода от простых стратегий агрегации (Average, Least Misery) к нейросетевым моделям, способным учитывать взаимодействия между участниками группы. Модели AGREE~\cite{cao2018agree}, GroupIM~\cite{sankar2020groupim} и AlignGroup~\cite{xu2024aligngroup} демонстрируют последовательное улучшение качества, при этом для сценария эфемерных групп~--- наиболее релевантного для настоящей работы~--- ключевую роль играет максимизация взаимной информации между представлениями пользователей и групп. Метрики справедливости~\cite{jin2023fairness_survey} должны быть включены в процедуру оценки наряду с традиционными метриками качества ранжирования.

В-пятых, проблема холодного старта решается несколькими взаимодополняющими способами: использованием контентных признаков (аудиоэмбеддинги в YAMBDA~\cite{yambda2025}), архитектурой DropoutNet~\cite{volkovs2017dropoutnet}, имитирующей отсутствие коллаборативных данных при обучении, и мета-обучением (MeLU~\cite{lee2019melu}), обеспечивающим быструю адаптацию к новым пользователям. Модели пространства состояний (Mamba4Rec~\cite{liu2024mamba4rec}) предлагают альтернативу трансформерам с линейной вычислительной сложностью, что критично при обработке длинных историй прослушиваний на ограниченных вычислительных ресурсах.

На основании проведённого обзора в настоящей работе принимается следующая архитектура системы групповых музыкальных рекомендаций:
\begin{enumerate}
    \item \textbf{Генерация кандидатов}~--- gSASRec с gBCE loss и механизмом персонализированной популярности (PPA), обученная на данных YAMBDA с использованием аудиоэмбеддингов для инициализации элементных представлений.
    \item \textbf{Ранжирование}~--- CatBoost с признаками взаимодействий, аудиоэмбеддингами и контекстными сигналами, обученный с учётом дисбаланса классов через \texttt{auto\_class\_weights}.
    \item \textbf{Групповая агрегация}~--- attention-based агрегация индивидуальных предпочтений с учётом метрик справедливости, обеспечивающая баланс между общим качеством рекомендаций и равномерностью удовлетворённости участников группы.
\end{enumerate}
